% !TeX program = xelatex
% 源：教学课件/A往年高一上学期课件/第4章/4.1.3 幂函数.pptx
\documentclass[aspectratio=169,10pt]{beamer}
\usepackage[UTF8,fontset=fandol]{ctex}
\IfFontExistsTF{Microsoft YaHei}{\setCJKsansfont{Microsoft YaHei}}{}
\usetheme{GaoZhongYuan}
% 教学内容层；原模板主题、图标及校园素材保持原样。
\usepackage{amsmath,amssymb,booktabs,array,pgfplots}
\pgfplotsset{compat=1.18}
\definecolor{lessonblue}{HTML}{225EA8}
\definecolor{lessonred}{HTML}{BA312D}
\definecolor{lessonpurple}{HTML}{82589F}
\setbeamerfont{normal text}{size=\large}
\setbeamerfont{block title}{size=\normalsize}
\setbeamerfont{block body}{size=\normalsize}
\setbeamertemplate{itemize item}{\color{campusgreen}\small$\bullet$}
\setbeamertemplate{enumerate item}{\color{campusgreen}\insertenumlabel.}
\setlength{\parskip}{.45em}
\setlength{\abovedisplayskip}{.5em}
\setlength{\belowdisplayskip}{.5em}
\newcommand{\R}{\mathbb R}
\newcommand{\Q}{\mathbb Q}
\newcommand{\ans}[1]{{\color{lessonred}#1}}
\newcommand{\key}[1]{{\color{campusgreen}\bfseries #1}}
\newcommand{\reason}[1]{{\color{lessonblue}#1}}
\newcommand{\gap}{\par\vspace{.6em}}
\pgfplotsset{lessonaxis/.style={axis lines=middle,axis line style={->},
 xlabel={$x$},ylabel={$y$},
 tick label style={font=\scriptsize},label style={font=\small},
 width=7.1cm,height=5.1cm,samples=150,clip=true,unbounded coords=jump,
 every axis plot/.append style={thick},
 legend style={font=\scriptsize,draw=none,fill=white,fill opacity=.9,text opacity=1},
 legend cell align=left}}
% 五个基本幂函数，均为可编辑矢量曲线。
\newcommand{\fivepowers}{%
\begin{tikzpicture}\begin{axis}[lessonaxis,xmin=-2.1,xmax=3.2,ymin=-2.1,ymax=3.6,
 xtick={-2,-1,1,2,3},ytick={-2,-1,1,2,3},legend pos=north west]
 \addplot[black,domain=-2:3.1]{x};\addlegendentry{$y=x$}
 \addplot[lessonblue,domain=-1.85:1.85]{x^2};\addlegendentry{$y=x^2$}
 \addplot[lessonred,domain=-1.3:1.5]{x^3};\addlegendentry{$y=x^3$}
 \addplot[campusgreen,domain=0:3.1]{sqrt(x)};\addlegendentry{$y=\sqrt{x}$}
 \addplot[lessonpurple,domain=.285:3.1]{1/x};\addlegendentry{$y=1/x$}
 \addplot[lessonpurple,domain=-2:-.48,forget plot]{1/x};
 \addplot[densely dashed,gray,forget plot] coordinates {(0,1)(1,1)(1,0)};
 \addplot[only marks,mark=*,mark size=1.3pt,forget plot] coordinates {(1,1)};
\end{axis}\end{tikzpicture}}
% 示意值只用于表达 0<q<m<1，未给原题补入数值条件。
\newcommand{\powerorder}{%
\begin{tikzpicture}\begin{axis}[lessonaxis,xmin=0,xmax=3.2,ymin=0,ymax=4,
 xtick={1,2,3},ytick={1,2,3},legend pos=north west]
 \addplot[lessonred,domain=0:1.95]{x^2};\addlegendentry{$y=x^p$}
 \addplot[campusgreen,domain=0:3.1]{x^.6};\addlegendentry{$y=x^m$}
 \addplot[lessonblue,domain=0:3.1]{x^.3};\addlegendentry{$y=x^q$}
 \addplot[lessonpurple,domain=.25:3.1]{1/x};\addlegendentry{$y=x^n$}
 \addplot[densely dashed,gray,forget plot] coordinates {(0,1)(1,1)(1,0)};
\end{axis}\end{tikzpicture}}
\newcommand{\expplot}[1][2]{%
\begin{tikzpicture}\begin{axis}[lessonaxis,xmin=-2.5,xmax=2.5,ymin=0,ymax=5,
 xtick={-2,-1,1,2},ytick={1,2,3,4},legend pos=north west]
 \addplot[lessonred,domain=-2.4:2.3]{2^x};\addlegendentry{$y=2^x$}
 \ifnum#1>1
 \addplot[lessonblue,domain=-2.3:2.4]{2^(-x)};\addlegendentry{$y=(1/2)^x$}
 \fi
 \ifnum#1>2
 \addplot[lessonpurple,domain=-2.4:1.45]{3^x};\addlegendentry{$y=3^x$}
 \addplot[campusgreen,domain=-1.45:2.4]{3^(-x)};\addlegendentry{$y=(1/3)^x$}
 \fi
 \addplot[only marks,mark=*,mark size=1.4pt,black,forget plot] coordinates {(0,1)};
\end{axis}\end{tikzpicture}}
% 四个选项依据原图特征重绘，B 为实立方根。
\newcommand{\rootchoices}{%
\foreach \letter/\kind in {A/1,B/2,C/3,D/4}{%
\begin{minipage}[t]{.24\textwidth}\centering
\begin{tikzpicture}\begin{axis}[lessonaxis,width=3.5cm,height=3.3cm,
 xmin=-1.6,xmax=1.7,ymin=-1.5,ymax=1.8,xtick={1},ytick={1},samples=201]
 \addplot[dashed,gray,domain=-1.4:1.6]{x};
 \ifnum\kind=1 \addplot[black,domain=-1.5:1.5]{tanh(x)};\fi
 \ifnum\kind=2 \addplot[black,domain=-1.5:1.5]{sign(x)*abs(x)^(1/3)};\fi
 \ifnum\kind=3 \addplot[black,domain=-1.12:1.2]{x^3};\fi
 \ifnum\kind=4 \addplot[black,domain=-1.5:1.5]{sign(x+.4)*abs(x+.4)^(1/3)};\fi
\end{axis}\end{tikzpicture}\\[-.5em]\letter
\end{minipage}}\par}

\title{4.1.3 幂函数}
\author{深圳中学高中园}
\date{}
\begin{document}
\campusmaketitle

% 原第2页
\begin{frame}[t]{任务一：生活中的幂函数}
 \renewcommand{\arraystretch}{1.65}
 \centering
 \begin{tabular}{lll}
  \toprule 情境 & 自变量 & 关系式\\\midrule
  铅笔单价为1元，购买 $x$ 支 & $x\in\mathbb N^*$ & $p=x$\\
  边长为 $x$ 的正方形 & $x>0$ & $S=x^2$\\
  棱长为 $x$ 的立方体 & $x>0$ & $V=x^3$\\\bottomrule
 \end{tabular}
 \pause
 \gap
 三个式子中，变量都在\key{底数}位置，指数是常数。
 \[y=x,\qquad y=x^2,\qquad y=x^3.\]
\end{frame}

% 原第3--4页，沿用原教材非零实数指数的口径。
\begin{frame}[t]{任务二：幂函数的定义}
 一般地，当 $\alpha$ 为非零实数时，函数
 \[\boxed{y=x^\alpha}\]
 叫做（$\alpha$ 次）幂函数。本课沿用此定义约定。
 \pause
 \begin{itemize}
  \item $x$ 是自变量，$\alpha$ 是常数。
  \item 系数为 $1$，底数为单个自变量 $x$。
  \item 只有 $x^\alpha$ 这一项，不另加常数项。
 \end{itemize}
 \[y=x^{-1},\quad y=x^{-2},\quad y=x^{1/2}=\sqrt x.\]
\end{frame}

% 原第5--6页
\begin{frame}[t]{任务三：五个基本幂函数}
 \begin{columns}[T,onlytextwidth]
  \begin{column}{.52\textwidth}\centering\fivepowers\end{column}
  \begin{column}{.45\textwidth}
   观察 $y=x,x^2,x^3,\sqrt x,1/x$：
   \begin{enumerate}
    \item 定义域和值域分别是什么？
    \item 哪些是奇函数、偶函数？
    \item 各自在哪些区间单调？
    \item 有哪些共同经过的点？
   \end{enumerate}
  \end{column}
 \end{columns}
\end{frame}

% 原第7页表格拆分，避免拥挤。
\begin{frame}[t]{任务三：定义域、值域与奇偶性}
 \centering\normalsize
 \renewcommand{\arraystretch}{1.65}
 \begin{tabular}{c|c|c|c}
  \toprule 函数 & 定义域 & 值域 & 奇偶性\\\midrule
  $y=x$ & $\R$ & $\R$ & 奇函数\\
  $y=x^2$ & $\R$ & $[0,+\infty)$ & 偶函数\\
  $y=x^3$ & $\R$ & $\R$ & 奇函数\\
  $y=\sqrt x$ & $[0,+\infty)$ & $[0,+\infty)$ & 非奇非偶\\
  $y=1/x$ & $(-\infty,0)\cup(0,+\infty)$ & $(-\infty,0)\cup(0,+\infty)$ & 奇函数\\\bottomrule
 \end{tabular}
 \gap
 \reason{判断奇偶性前，先检查定义域是否关于原点对称。}
\end{frame}

\begin{frame}[t]{任务三：单调区间与定点}
 \normalsize\renewcommand{\arraystretch}{1.7}
 \centering
 \begin{tabular}{c|l|c}
  \toprule 函数 & 单调性 & 经过的点\\\midrule
  $y=x$ & 在 $\R$ 上递增 & $(0,0),(1,1)$\\
  $y=x^2$ & 在 $(-\infty,0]$ 上递减，在 $[0,+\infty)$ 上递增 & $(0,0),(1,1)$\\
  $y=x^3$ & 在 $\R$ 上递增 & $(0,0),(1,1)$\\
  $y=\sqrt x$ & 在 $[0,+\infty)$ 上递增 & $(0,0),(1,1)$\\
  $y=1/x$ & 在 $(-\infty,0)$、$(0,+\infty)$ 上分别递减 & $(1,1)$\\\bottomrule
 \end{tabular}
 \gap
 \ans{$1/x$ 的两个单调区间不能合写成一个单调区间。}
\end{frame}

% 原第8、13页；纠正把所有负整数次幂都按奇函数描述的问题。
\begin{frame}[t]{任务四：整数次幂的奇偶性}
 对非零整数 $k$，$y=x^k$ 满足
 \[f(-x)=(-x)^k=(-1)^k x^k.\]
 \pause
 \begin{itemize}
  \item $k$ 为奇数：奇函数，图象关于原点对称。
  \item $k$ 为偶数：偶函数，图象关于 $y$ 轴对称。
 \end{itemize}
 \gap
 当 $k<0$ 时，$x=0$ 不在定义域内；两支曲线分别趋近坐标轴。\\
 \ans{$y=x^{-1}$ 的左支在第三象限，$y=x^{-2}$ 的左支在第二象限。}
\end{frame}

% 原第9、11、13页的重复内容合并，明确只是非负/正半轴性质。
\begin{frame}[t]{任务五：正半轴上的性质}
 对 $y=x^\alpha\ (\alpha\ne0)$：
 \begin{itemize}
  \item 当 $\alpha>0$ 时，在 $[0,+\infty)$ 上有定义且递增，\\
        这一部分的值域为 $[0,+\infty)$，过 $(0,0),(1,1)$。
  \item 当 $\alpha<0$ 时，在 $(0,+\infty)$ 上有定义且递减，\\
        这一部分的值域为 $(0,+\infty)$，过 $(1,1)$。
 \end{itemize}
 \pause
 \gap
 完整定义域、完整值域还需结合具体指数判断。\\
 例如 $y=x^3$ 的值域为 $\R$，$y=x^2$ 的值域为 $[0,+\infty)$。
\end{frame}

% 原第10--11页
\begin{frame}[t]{任务五：象限与定义域}
 对任意允许的指数 $\alpha$，当 $x>0$ 时，$x^\alpha>0$。
 \pause
 \begin{itemize}
  \item 第一象限一定有图象，第四象限一定没有图象。
  \item $x<0$ 时是否有定义，以及函数值的符号，取决于指数。
  \item 分数次幂按相应根式判断，例如
   \[x^{1/2}=\sqrt x\ (x\ge0),\qquad x^{1/3}=\sqrt[3]x\ (x\in\R).\]
 \end{itemize}
 \reason{不能把第一象限中的性质直接当成整个定义域上的性质。}
\end{frame}

% 原第12--14页
\begin{frame}[t]{任务六：指数与图象高低}
 对同一个 $x>0$，设 $\alpha<\beta$。
 \pause
 \begin{columns}[T,onlytextwidth]
  \begin{column}{.47\textwidth}
   \key{当 $0<x<1$ 时}
   \[\ans{x^\alpha>x^\beta}.\]
   指数越大，图象越低。
  \end{column}
  \begin{column}{.47\textwidth}
   \key{当 $x>1$ 时}
   \[\ans{x^\alpha<x^\beta}.\]
   指数越大，图象越高。
  \end{column}
 \end{columns}
 \gap
 当 $x=1$ 时，$1^\alpha=1^\beta=1$。\\
 例如 $0<x<1$ 时 $x^3<x^2<x<\sqrt x$；$x>1$ 时次序反向。
\end{frame}

% 原第13页的代数解释。
\begin{frame}[t]{任务六：同指数幂的比较依据}
 设 $a>b>0$，则 $a/b>1$。
 \pause
 \begin{align*}
  r>0&:\quad \frac{a^r}{b^r}=\left(\frac ab\right)^r>1
       \quad\Longrightarrow\quad a^r>b^r,\\[.8em]
  r<0&:\quad \frac{a^r}{b^r}=\left(\frac ab\right)^r<1
       \quad\Longrightarrow\quad a^r<b^r.
 \end{align*}
 \gap
 \reason{这说明了正半轴上幂函数的增减性。}
\end{frame}

% 原第15页为第二课时封面。
\begin{frame}[t]{第二课时：幂函数的应用}
 \begin{enumerate}
  \item 由定义识别幂函数，求参数。
  \item 根据图象判断指数，利用已知点求解析式。
  \item 利用单调性比较大小、解不等式。
 \end{enumerate}
 \gap
 \key{每次使用性质前，先确定定义域与单调区间。}
\end{frame}

% 原第16页(1)
\begin{frame}[t]{例1：幂函数的概念}
 （多选）下列函数中是幂函数的是（\quad）。
 \begin{align*}
  \mathrm A.\ y&=\left(\frac12\right)^x;&
  \mathrm B.\ y&=x^5;\\[.6em]
  \mathrm C.\ y&=4x^2;&\mathrm D.\ y&=x.
 \end{align*}
 \pause
 \ans{答案：B、D。}
 \gap
 A 中变量位于指数位置；C 中 $x^2$ 的系数为 $4$。\\
 B、D 均满足本课定义的标准形式 $y=x^\alpha$。
\end{frame}

% 原第16页(2)
\begin{frame}[t]{例1：由定义求参数}
 已知
 \[y=(m^2+2m-2)x^{m^2-2}+2n-3\]
 是幂函数，求 $m,n$ 的值。
 \pause
 由系数为 $1$、常数项为 $0$，得
 \[\begin{cases}m^2+2m-2=1,\\2n-3=0.\end{cases}\]
 所以 $\ans{m=-3\text{ 或 }1,\quad n=\frac32}$。\\
 两种情况下指数分别为 $7,-1$，均非零，符合本课定义。
\end{frame}

% 原第17页
\begin{frame}[t]{跟踪训练1：定义与单调性}
 若 $y=(m^2+2m-2)x^m$ 为幂函数，且在第一象限为增函数，\\
 则 $m$ 的值为（\quad）。
 \[\mathrm A.\ 1\qquad\mathrm B.\ {-3}\qquad
   \mathrm C.\ {-1}\qquad\mathrm D.\ 3\]
 \pause
 \[\begin{cases}m^2+2m-2=1,\\m>0,\end{cases}
   \quad\Longrightarrow\quad\ans{m=1}.\]
 \ans{答案：A。}
\end{frame}

% 原幂函数第18--20、23、25--27页；也用于指数函数原稿第2--5页复习。
\begin{frame}[t]{例2：立方根函数的图象}
 函数 $y=x^{1/3}=\sqrt[3]x$ 的图象是（\quad）。
 \par\rootchoices
 \pause
 \key{过 $(0,0)$、$(1,1)$，关于原点对称。}\\
 当 $x>1$ 时，$x^{1/3}<x$；当 $0<x<1$ 时，$x^{1/3}>x$。\\
 因而应选 \ans{B}。
\end{frame}

\begin{frame}[t]{例2：由图象判断指数大小}
 \begin{columns}[T,onlytextwidth]
  \begin{column}{.53\textwidth}\centering\powerorder\end{column}
  \begin{column}{.44\textwidth}
   根据图象中 $x>1$ 的部分，比较 $m,n,p,q$ 的大小。
   \pause
   \[n<0<q<m<1<p.\]
   因而
   \[\ans{n<q<m<p}.\]
   \reason{在 $x>1$ 时，指数越大，图象越高。}
  \end{column}
 \end{columns}
\end{frame}

\begin{frame}[t]{例3：待定指数求函数值}
 若 $f(x)$ 是幂函数，且 $f(8)=16$，求 $f(9),f(64)$。
 \pause
 设 $f(x)=x^\alpha$。由 $8^\alpha=16$ 得
 \[2^{3\alpha}=2^4\quad\Longrightarrow\quad\alpha=\frac43.\]
 因此 $f(x)=x^{4/3}$，
 \[\ans{f(9)=9^{4/3}=9\sqrt[3]9},\qquad
   \ans{f(64)=(2^6)^{4/3}=256}.\]
 \reason{函数类型已知时，先设标准形式，再求参数。}
\end{frame}

\begin{frame}[t]{例4：同指数幂的大小}
 比较各组中两个数的大小：
 \[(1)\ 1.5^{1.4},\ 1.6^{1.4};\qquad
   (2)\ 1.5^{0.4},\ 1.6^{0.4}.\]
 \pause
 \begin{itemize}
  \item $y=x^{1.4}$ 在 $[0,+\infty)$ 上递增，故
   \[\ans{1.5^{1.4}<1.6^{1.4}}.\]
  \item $y=x^{0.4}$ 在 $[0,+\infty)$ 上递增，故
   \[\ans{1.5^{0.4}<1.6^{0.4}}.\]
 \end{itemize}
\end{frame}

\begin{frame}[t]{例4：负指数与比较方法}
 \[(3)\ 1.5^{-1.5},\ 1.6^{-1.5}.\]
 \pause
 $y=x^{-1.5}$ 在 $(0,+\infty)$ 上递减，且 $1.5<1.6$，故
 \[\ans{1.5^{-1.5}>1.6^{-1.5}}.\]
 \gap
 \key{用单调性直接比较时，两个自变量应处于同一单调区间。}
 \gap
 跨区间时应另作分析，例如：$-1<1$，但 $(-1)^{-1}<1^{-1}$。\\
 $y=1/x$ 虽在两侧分别递减，却不在整个定义域上递减。
\end{frame}


% 原第21--22页，顺序调到典例之后。
\begin{frame}[t]{跟踪训练2：图象与指数}
 \begin{columns}[T,onlytextwidth]
  \begin{column}{.5\textwidth}
   $y=x^n$ 中 $n$ 取 $2,\frac12,-2$。\\
   曲线 $c_1,c_2,c_3$ 的指数依次为（\quad）。
   \normalsize
   \begin{align*}
    \mathrm A.&\ {-2},\tfrac12,2;&\mathrm B.&\ 2,\tfrac12,-2;\\
    \mathrm C.&\ {-2},2,\tfrac12;&\mathrm D.&\ 2,-2,\tfrac12.
   \end{align*}
   \pause
   取 $x=4$，三个值为 $16,2,\frac1{16}$。\\
   \ans{答案：B。}
  \end{column}
  \begin{column}{.48\textwidth}
   \begin{tikzpicture}\begin{axis}[lessonaxis,xmin=0,xmax=3.2,ymin=0,ymax=4,
    xtick={1},ytick={1},legend pos=north west]
    \addplot[lessonred,domain=0:1.98]{x^2};\addlegendentry{$c_1$}
    \addplot[campusgreen,domain=0:3.1]{sqrt(x)};\addlegendentry{$c_2$}
    \addplot[lessonblue,domain=.51:3.1]{x^(-2)};\addlegendentry{$c_3$}
    \addplot[dashed,gray,forget plot] coordinates {(0,1)(1,1)(1,0)};
   \end{axis}\end{tikzpicture}
  \end{column}
 \end{columns}
\end{frame}

\begin{frame}[t]{跟踪训练2：图象所在象限}
 当 $\alpha\in\{-1,\frac12,1,2,3\}$ 时，幂函数 $y=x^\alpha$ 的图象\\
 不可能经过第\underline{\hspace{1cm}}象限。
 \pause
 \begin{itemize}
  \item $y=x^{-1},x,x^3$ 经过第一、三象限。
  \item $y=\sqrt x$ 经过第一象限。
  \item $y=x^2$ 经过第一、二象限。
 \end{itemize}
 \ans{答案：四。}
\end{frame}

% 原第24页
\begin{frame}[t]{跟踪训练3：由点求值}
 已知幂函数 $f(x)$ 的图象经过点 $\left(3,\frac19\right)$，求 $f(4)$。
 \pause
 设 $f(x)=x^\alpha$，则
 \[3^\alpha=\frac19=3^{-2}\quad\Longrightarrow\quad\alpha=-2.\]
 因而
 \[\ans{f(4)=4^{-2}=\frac1{16}}.\]
\end{frame}

% 原第27--28页
\begin{frame}[t]{跟踪训练4：借助中间量比较}
 将下列各数从小到大排列：
 \[\left(\frac23\right)^{-1/3},\quad
   \left(\frac35\right)^{1/2},\quad
   \left(\frac25\right)^{1/2},\quad
   \left(\frac76\right)^0.\]
 \pause
 $\left(\frac76\right)^0=1$，$\left(\frac23\right)^{-1/3}>1$；另外
 \[0<\left(\frac25\right)^{1/2}<\left(\frac35\right)^{1/2}<1.\]
 因此
 \[\ans{\left(\frac25\right)^{1/2}<\left(\frac35\right)^{1/2}
 <\left(\frac76\right)^0<\left(\frac23\right)^{-1/3}}.\]
\end{frame}

% 原第29页；用完整符号表避免漏分情况。
\begin{frame}[t]{例5：倒数不等式}
 已知 $(a+1)^{-1}<(3-2a)^{-1}$，求 $a$ 的取值范围。
 \pause
 先排除 $a=-1,\frac32$，移项得
 \[\frac1{a+1}-\frac1{3-2a}
   =\frac{2-3a}{(a+1)(3-2a)}<0.\]
 \renewcommand{\arraystretch}{1.4}\centering\normalsize
 \begin{tabular}{c|cccc}
  $a$ & $(-\infty,-1)$ & $(-1,\frac23)$ & $(\frac23,\frac32)$ & $(\frac32,+\infty)$\\\hline
  分式符号 & $-$ & $+$ & $-$ & $+$
 \end{tabular}
 \[\ans{a\in(-\infty,-1)\cup\left(\frac23,\frac32\right)}.\]
\end{frame}

% 原第30页
\begin{frame}[t]{解不等式的方法}
 \begin{enumerate}
  \item 先由根式、分母等确定有意义的条件。
  \item 确认两个自变量是否处于同一单调区间。
  \item 同一区间内用单调性转化；跨区间时分类或比较符号。
  \item 将所得范围与有意义的条件取交集，检查端点。
 \end{enumerate}
 \gap
 \reason{例如倒数函数两侧分别递减，但不能直接跨越零点反向比较。}
\end{frame}

% 原第31页；纠正解答中把非负条件写成严格大于零。
\begin{frame}[t]{跟踪训练5：根式不等式}
 若 $(3-2m)^{1/2}>(m+1)^{1/2}$，求实数 $m$ 的取值范围。
 \pause
 $y=\sqrt x$ 在 $[0,+\infty)$ 上递增，所以原不等式等价于
 \[\begin{cases}3-2m\ge0,\\m+1\ge0,\\3-2m>m+1.\end{cases}\]
 解得
 \[\ans{m\in\left[-1,\frac23\right)}.\]
\end{frame}

\begin{frame}[t]{课堂小结}
 \begin{itemize}
  \item 幂函数：变量在底数位置，先核对系数与附加项。
  \item 定义域与奇偶性：结合具体指数判断。
  \item 正半轴单调性：正指数递增，负指数递减。
  \item 比较大小：同指数用幂函数，必要时借助 $1$ 等中间量。
  \item 解不等式：先定范围，再用性质，最后检查端点。
 \end{itemize}
\end{frame}
\campusclosing[谢谢]
\end{document}
